\clearpage
\section*{历史注记}
\phantomsection
\addcontentsline{toc}{section}{第 5 节历史注记}

（第三章，第5节）

（方括号中的数字指本注记末尾的参考文献。）

与整数和基数概念相关的思想演变，与集合论和数理逻辑的历史密不可分，读者可参阅第四章后的历史注记。本注记的目的在于简要指出记数法和“组合分析”历史中的一些突出事实。

历史和考古向我们揭示了大量的“记数系统”，其主要目的是为每个整数（直至某个取决于实际使用需求的界限）赋予一个名称和一种书面表示，这种表示由有限数量的符号按照或多或少规则的法则构成。最为常见的做法是将整数分解为“逐级单位”的和 \(b_{1}, b_{2}, \ldots, b_{n}, \ldots\) ，其中每一项都是其前一项的整数倍；尽管通常 \(b_{n}/b_{n-1}\) 取为固定的数 b（系统的“基数”，通常为10），但已知有许多违反此规则的例外。例如，在巴比伦系统中， \(b_{n}/b_{n-1}\) 有时为10，有时为6 {[}1{]}，而在玛雅人的年代学系统中， \(b_{n}/b_{n-1}\) 除 n = 2 外均等于20，且 \(b_{2}/b_{1} = 18\) {[}2{]}。至于相应的书写符号，它必须指明每一阶 i 的“单位” \(b_{i}\) 的数目。在许多系统中（例如埃及、希腊和罗马系统），逐级倍数 \(k.b_{i}\) （其中 k 从1变化到 \((b_{i+1}/b_{i}) - 1\) ）由同时依赖于 k 和 i 的符号表示。第一个重要进步是用同一符号表示所有数 \(k.b_{i}\) （对于相同的 k 值）：这就是“按位记数”的原理，其中阶指标 i 由表示 \(k.b_{i}\) 的符号出现在“第 i 位”这一事实来指示。第一个此类系统是巴比伦人的系统，他们至少在公元前2000年就用同一符号表示所有倍数 \(k.60^{\pm i}\) 对应于指数 i 的不同取值（{[}1{]}，第 93-109 页）。这种体系的缺点当然在于其模糊性，只要没有任何标记来指示某一阶的单位是否存在，即只要该体系没有通过引入“零”来完善，这种模糊性就存在。尽管如此，巴比伦人在其历史的大部分时间里都没有这样的记号，并且在公元前最后两个世纪之前都没有使用“零”，而且即使在那时，也仅在一个数的内部使用；直到那时，只有上下文才能阐明所讨论符号的含义。只有另外两种体系系统地使用了“零”：玛雅人的体系（显然自公元纪年开始就一直在使用 {[}2{]}）以及我们现在的十进制体系，后者（经由阿拉伯人）源自印度数学，在那里，零的使用自公元最初几个世纪起就有确证。此外，将零视为一个数（而不仅仅是一个分隔符号）的概念以及将其引入计算，是印度人的原创贡献 {[}3{]}。当然，一旦“按位记数”的原理被掌握，将其推广到任意基数就很容易了。自 17 世纪以来提出的各种“基数”优劣的讨论，取决于数值计算的技术，此处无法深入探讨。我们仅指出，这些体系根基所在的运算，即所谓的“欧几里得除法”，直到希腊时代才出现，并且无疑可以追溯到早期毕达哥拉斯学派，他们将其作为其理论算术的基本工具。

一般性的计数问题，统称为“组合分析”，似乎直到古典时代的最后几个世纪才有人尝试；只有公式 \(\binom{n}{2} = \frac{1}{2} n(n - 1)\) 在公元三世纪得到证实。印度数学家婆什迦罗（12 世纪）知道一般公式 \(\binom{n}{p}\) 。在利未·本·格尔森的一份手稿（13 世纪初）中发现了更系统的研究：他得到了从 n 个对象中每次取 p 个的排列数 \(V_{n}^{p}\) 的归纳公式，特别是 n 个对象的排列数，并且他陈述了与关系 \(\binom{n}{p} = V_{n}^{p}/p!\) 和 \(\binom{n}{p} = \binom{n}{n-p}\) 等价的规则 ({[}4{]}，第 64-65 页）。但这份手稿似乎并未被其同时代人所知，这些结果在后来的几个世纪里才被数学家们逐渐重新发现。至于后来的进展，让我们记录一下卡尔达诺证明了 n 个元素的集合的非空子集的数量是 \(2^{n} - 1\) 。帕斯卡和费马在创立概率论时，重新发现了 \(\binom{n}{p}\) 的表达式，而帕斯卡是第一个观察到这些数与二项式定理之间关系的人，这个定理似乎自 13 世纪起就为阿拉伯人所知，14 世纪为中国人所知，并在 16 世纪初在西方被重新发现，同时还有被称为“帕斯卡三角”的系数的归纳计算方法（{[}4{]}，第 35-38 页）。最后，大约在 1676 年，莱布尼茨得到了（但未发表）“多项式系数”的一般公式，该公式在 20 年后由棣莫弗独立重新发现并发表。

\section*{参考书目}
\phantomsection
\addcontentsline{toc}{section}{参考书目}

\begin{enumerate}
\def\labelenumi{\arabic{enumi}.}
\item
  O. 诺伊格鲍尔，《古代数学史讲义》，第 I 卷：前希腊数学，柏林（施普林格），1934 年。
\item
  S. G. 莫利，《古代玛雅》，斯坦福大学出版社，1946 年。
\item
  B. 达塔 和 A. N. 辛格，《印度数学史》，第 I 卷，拉合尔（莫蒂拉尔·巴纳西达斯），1935 年。
\item
  J. 特罗普克，《初等数学史》，第 VI 卷：分析，解析几何，柏林-莱比锡（德·格鲁伊特），1924 年。
\end{enumerate}
